% !TEX root = ../main.tex
\chapter{等式を命題として読む}
\label{chap:ch03_equations_as_props}
\BookIndex{とうしき}{等式}
\BookIndex{めいだい}{命題}

\begin{chapterabstract}
本章では、\texttt{hasDiscount true = true} のような\emph{等式}を読む練習をします。
プログラミングの多くの場面で \texttt{=} は代入を意味しますが、Lean で仕様を書くときの \texttt{=} は代入ではありません。
「左辺と右辺が等しい」という\emph{主張}です。
本章では、この主張を「命題」と呼び、成立する主張と成立しない主張を式の展開によって区別します。
\end{chapterabstract}

\begin{learninggoals}
\item \texttt{Bool} を返す関数の結果を \texttt{\#eval} で確認できます。
\item \texttt{hasDiscount true = true} を「等しいという主張」として読めます。
\item \texttt{=} が代入ではなく、成り立つかどうかを問える主張であると説明できます。
\end{learninggoals}

\section{関数の結果を等式で主張する}

仕様で確認したい内容は、しばしば「ある入力でこの関数はこの値になる」という形の主張です。
たとえば、「会員なら割引対象になる」は、\texttt{hasDiscount true} が \texttt{true} になると言い換えられます。
この「\texttt{hasDiscount true} は \texttt{true} と等しい」を、Lean では \texttt{hasDiscount true = true} と書きます。
本章では、この一行を正しく読めるようにします。

図\ref{fig:ch03-overview}は、関数を評価して値を得ることと、その値について等式を主張することを分けて示します。

\FigurePlaceholder{fig-ch03-overview.png}{0.90}{値・計算・命題}{fig:ch03-overview}

\section{まず関数と実行結果を見る}

題材には、会員かどうかで割引対象を判定する小さな関数を使います。
ここでは、引数 \texttt{member} をそのまま返す単純な関数にします。

\begin{listing}[htbp]
\begin{minted}{lean4}
def hasDiscount (member : Bool) : Bool :=
  member

#eval hasDiscount true
#eval hasDiscount false
\end{minted}
\caption{\texttt{hasDiscount}}
\label{lst:ch03_hasdiscount}
\end{listing}

コード\ref{lst:ch03_hasdiscount}の \texttt{hasDiscount} は、\texttt{member} を受け取り、そのまま返す関数です。
したがって、\texttt{hasDiscount true} は \texttt{true} に、\texttt{hasDiscount false} は \texttt{false} に計算されます。
ここでの \texttt{member} は、呼び出し時に \texttt{true} または \texttt{false} を受け取る引数です。

\texttt{\#eval hasDiscount true} は、関数を実際に\emph{計算}して値を確認します。
これは、テストで特定の入力の結果を確かめる感覚に近い操作です。
この計算結果を「主張」として書き直します。

\section{等式を主張として読む}

\texttt{hasDiscount true} が \texttt{true} になることを、Lean では次のような\emph{主張}として書けます。

\begin{listing}[htbp]
\begin{minted}{lean4}
def hasDiscount (member : Bool) : Bool :=
  member

#check (hasDiscount true = true)
\end{minted}
\caption{\texttt{hasDiscount true = true}}
\label{lst:ch03_equation}
\end{listing}

コード\ref{lst:ch03_equation}の \texttt{hasDiscount true = true} は、代入ではありません。
「\texttt{hasDiscount true} の結果は \texttt{true} と等しい」という主張です。
\texttt{\#check} で確認すると、Lean がこれを値ではなく、成り立つかどうかを問える\emph{命題}として扱っていることが分かります。

この一行は、次のように読めます。
段階を分けて読むと、この主張が成り立つ理由を確認できます。

\begin{table}[htbp]
\centering
\caption{\texttt{hasDiscount true = true}}
\label{tab:ch03-reading}
\begin{tabular}{ll}
\toprule
読む段階 & Lean が見ている形 \\
\midrule
最初の主張 & \texttt{hasDiscount true = true} \\
\texttt{hasDiscount} を展開 & \texttt{true = true} \\
結論 & 左右が同じなので成り立つ \\
\bottomrule
\end{tabular}
\end{table}

\section{成り立たない主張も書ける}

主張を書けることと、その主張が成り立つことは別です。
たとえば、非会員でも割引があるという誤った主張も、式としては書けます。

\begin{table}[htbp]
\centering
\caption{\texttt{hasDiscount true = true} と \texttt{hasDiscount false = true}}
\label{tab:ch03-true-false}
\begin{tabular}{lll}
\toprule
書いた主張 & 展開後 & 成り立つか \\
\midrule
\texttt{hasDiscount true = true} & \texttt{true = true} & 成り立つ \\
\texttt{hasDiscount false = true} & \texttt{false = true} & 成り立たない \\
\bottomrule
\end{tabular}
\end{table}

表\ref{tab:ch03-true-false}のように、\texttt{hasDiscount false = true} は展開すると \texttt{false = true} になり、左右が異なるため成り立ちません。
主張はまず正しく読める必要があり、定義を展開して左右を比較すると成否を判断できます。

\section{本章のまとめ}

要点と記法
\begin{itemize}
\item \texttt{=}：代入ではなく、左右が等しいという主張
\item \texttt{\#check (式 = 式)}：等式が命題として扱われることの確認
\item 定義を展開した左右の確認：書いた命題が成り立つかを判断する方法
\end{itemize}
